\documentclass[11pt]{article}
\usepackage[margin=0.9in]{geometry}
\usepackage{amsmath,amssymb,amsthm,mathtools,hyperref}
\hypersetup{hidelinks,pdftitle={Q222: exact canonical completion and polynomial-selector obstruction}}
\newtheorem{theorem}{Theorem}
\newtheorem{proposition}[theorem]{Proposition}
\newtheorem{corollary}[theorem]{Corollary}
\newcommand{\PSD}{\succeq0}
\newcommand{\R}{\mathbb R}
\title{An exact two-LMI reduction of the canonical $Q_{222}$ slack\\
and an obstruction to every polynomial Bob selector}
\author{}
\date{}
\begin{document}
\maketitle
\vspace{-2em}
\begin{abstract}
A source-independent polynomial state change reduces the complete canonical
slack, with all nine source coefficients and four independent state coordinates,
to a matrix quadratic in the squared Alice opening.  Positivity in Alice is
equivalent to two $4\times4$ matrix inequalities with one $2\times2$ Bob-dependent
selector.  This reduction includes all zero source weights and all projective
testing endpoints.  On C1, every polynomial selector reduces to three scalar
parameters.  An already certified rational valid source admits none: its exact
contact fixes those parameters, and a $3\times3$ principal minor is negative at
$v=1$.  Thus polynomial completion of this normal form fails; a fixed algebraic
selector or an additional fixed testing multiplier remains unresolved.
\end{abstract}

\section{The full state reduction}
Let $c=(d,\lambda,n,m,C,A,D,\alpha,\beta)\geq0$ and put
$D_u=1+u^2$, $D_v=1+v^2$.  The established complete canonical parent is
\begin{align}
F_c={}&4\lambda D_u(W+vX)^2+4nD_uD_vX^2
 +4m(vW-X+uvY-uZ)^2\notag\\
&+4CD_v(uW-Y)^2+4AD_uD_vY^2
 +(4D+2d)(uW+uvX-Y-vZ)^2\notag\\
&+(4\alpha+2d)D_u(vY-Z)^2
 +(4\beta+2d)D_v(X+uZ)^2-2dD_uD_vZ^2.\label{parent}
\end{align}
Define the symmetric
$2\times2$ matrices
\begin{align}
P(v)&=\begin{pmatrix}
4mv^2&-4mv\\-4mv&4m+(4\beta+2d)D_v
\end{pmatrix},\notag\\
N(v)&=\begin{pmatrix}
4CD_v+4D+2d&(4D+2d)v\\(4D+2d)v&(4D+2d)v^2
\end{pmatrix},\label{blocks}\\
E(v)&=\begin{pmatrix}
4AD_v+(4\alpha+2d)v^2&-(4\alpha+2d)v\\
-(4\alpha+2d)v&4\alpha-2dv^2
\end{pmatrix},\notag\\
K(v)&=4\lambda\begin{pmatrix}1&v\\v&v^2\end{pmatrix}
       +4nD_v\begin{pmatrix}0&0\\0&1\end{pmatrix}.\notag
\end{align}
All entries are affine in $c$ and quadratic in the testing variable $v$.
For $a=(S,T)^T$, $b=(P_0,Q_0)^T$, and $U=u^2$, make the polynomial change
\begin{equation}
W=S+UP_0,\quad X=T+UQ_0,\quad
Y=u(S-P_0),\quad Z=u(T-Q_0).\label{change}
\end{equation}
The state coordinates $P_0,Q_0$ are unrelated to the matrix $P(v)$.
Direct substitution gives the exact identity
\begin{equation}
F_c=D_u\Phi_c(U,v;a,b),\qquad
\boxed{\begin{aligned}
\Phi_c={}&(1+U)a^TPa+U(1+U)b^TNb\\
 &+U(a-b)^TE(a-b)+(a+Ub)^TK(a+Ub).
\end{aligned}}\label{normal}
\end{equation}
For example, the $m,C,D,\alpha,\beta$ event roots in \eqref{parent}
become respectively
$D_u(vS-T)$, $uD_uP_0$, $uD_u(P_0+vQ_0)$,
$u\{v(S-P_0)-(T-Q_0)\}$, and $D_uT$.
These substitutions prove \eqref{normal} without a source-dependent inverse.

For completeness, write $\psi=(W,X,Y,Z)^T$ and let $\mathcal P_c(U,v)$
be the symmetric matrix of $\Phi_c$ in $(a,b)$. Its three coefficients are
\begin{equation}
\mathcal P_c=
\begin{pmatrix}P+K&0\\0&0\end{pmatrix}
+U\begin{pmatrix}P+E&K-E\\K-E&N+E\end{pmatrix}
+U^2\begin{pmatrix}0&0\\0&N+K\end{pmatrix}.\label{quadratic}
\end{equation}
The change matrix has determinant $u^2D_u^2$.  Its polynomial inverse
numerator is
\begin{equation}
L(u)=\begin{pmatrix}
u&0&u^2&0\\0&u&0&u^2\\u&0&-1&0\\0&u&0&-1
\end{pmatrix},\qquad
\boxed{u^2D_uF_c(\psi)=(L(u)\psi)^T
\mathcal P_c(u^2,v)(L(u)\psi).}\label{multiplier}
\end{equation}
Equation \eqref{multiplier} is a polynomial normal-form identity. It is not a
sum-of-squares certificate.

\section{Two matrix inequalities, with the quantifiers retained}
\begin{theorem}\label{thm:selector}
For any fixed real source $c$ and fixed real Bob test $v$, the following are
equivalent:
\begin{enumerate}
\item $\mathcal P_c(U,v)\PSD$ for every $U\geq0$.
\item There exists a real $2\times2$ matrix $J$ such that
\begin{equation}
\boxed{\mathcal G_J=
\begin{pmatrix}P+K&J\\J^T&N+K\end{pmatrix}\PSD,
\qquad
\mathcal H_J=
\begin{pmatrix}P+E&K-E-J\\K-E-J^T&N+E\end{pmatrix}\PSD.}\label{lmis}
\end{equation}
\end{enumerate}
The equivalence also holds over every real closed field.  In particular,
complete canonical validity is exactly
\begin{equation}
\boxed{\forall v\in\R\ \exists J\in\R^{2\times2}:
\mathcal G_J(v)\PSD,\quad\mathcal H_J(v)\PSD.}\label{quantifiers}
\end{equation}
All source faces, including $d=0$ and arbitrary zero event weights, are retained.
\end{theorem}
\begin{proof}
The reverse implication is the explicit identity
\begin{equation}
\Phi_c=
\binom a{Ub}^{\!T}\mathcal G_J\binom a{Ub}
+U\binom ab^{\!T}\mathcal H_J\binom ab.\label{split}
\end{equation}
For the forward implication, $\mathcal P_c(t^2,v)$ is a positive semidefinite
univariate matrix polynomial of degree at most four.  The real-line matrix
Fej\'er--Riesz theorem gives a polynomial matrix factor of degree at most two;
real and imaginary parts give a real sum of matrix squares
\cite[Theorem 1.2]{Zalar}. Averaging its values at $t$ and $-t$ separates even
and odd powers. Consequently
\[
\mathcal P_c(U,v)=(B_0+UB_2)^T(B_0+UB_2)+UB_1^TB_1
\]
for rectangular real matrices $B_0,B_1,B_2$.
The constant coefficient in \eqref{quadratic} has its last two rows and columns
zero, so $B_0=[R\ \ 0]$.  The coefficient of $U^2$ has its first two rows and
columns zero, so $B_2=[0\ \ T]$.  This uses only the elementary fact that a
zero diagonal entry of a positive semidefinite Gram matrix forces its entire
row and column to vanish.  Setting $J=R^TT$ gives
$\mathcal G_J=[R\ T]^T[R\ T]$ and
$\mathcal H_J=B_1^TB_1$, proving \eqref{lmis}.

This is a fixed finite-dimensional first-order assertion in the real closed
field language: positivity is expressed by principal minors, and the
half-line condition by quantification over $U$ and a four-vector.  Transfer
therefore proves precisely the same assertion over every real closed field.

For $u\ne0$, \eqref{change} is invertible and so positivity of $F_c$ is equivalent
to positivity of $\mathcal P_c(u^2,v)$.  The excluded value $u=0$ is recovered
by polynomial continuity; all $U>0$ are attained.  More explicitly,
\eqref{parent} has bidegree at most $(2,2)$ in $(u,v)$, so
\[
F_c^{\mathrm h}(r,s;\psi)
=r_0^2s_0^2F_c(r_1/r_0,s_1/s_0;\psi)
\]
is polynomial in the homogeneous angle pairs $r=(r_0,r_1)$ and
$s=(s_0,s_1)$.  On their unit circles it equals the normalization
$F_c/(D_uD_v)$.  Continuity recovers the Alice and Bob projective endpoints.
Over a real closed field the same endpoint statement follows by the sign of
the first nonzero coefficient in the corresponding local polynomial chart.
No source coefficient was divided out.
\end{proof}

There is also an elementary check on the face $d=0$: then $E\PSD$, and
$J=K$ makes both matrices in \eqref{lmis} sums of the positive diagonal blocks
$P,N$ and a square block built from $K$ or $E$. More generally adding a
nonnegative event source with $d=0$ preserves any certificate by adding its
$K$ to $J$.  The unresolved completion is therefore associated with the signed
core, not with adding positive events.

Theorem~\ref{thm:selector} does not choose $J$ in a fixed space of functions
of $v$. Its real-closed-field proof permits coefficients algebraic in both
source and testing data. The tensor product of the independent source and
testing fields is not a real closed field, and this proof supplies no
sum-of-squares factorization in that tensor product.

\section{The complete polynomial-selector obstruction}
\begin{proposition}\label{prop:degree}
If polynomial entries $J(v)$ satisfy \eqref{lmis} for every real $v$, they have
degree at most two.  A selector can then be chosen of the form
\[
J(v)=\begin{pmatrix}p_0+p_2v^2&q_1v\\r_1v&s_0+s_2v^2\end{pmatrix}.
\]
On the entire C1 face $\lambda=n=0$, existence of any polynomial selector is
equivalent to existence of one with just three scalar parameters:
\begin{equation}
\boxed{J(v)=\begin{pmatrix}pv^2&0\\rv&sv^2\end{pmatrix}.}\label{three}
\end{equation}
\end{proposition}
\begin{proof}
Every $2\times2$ principal minor of $\mathcal G_J$ gives
$|J_{ij}(v)|^2\leq(P+K)_{ii}(N+K)_{jj}$.  The right side has degree at most
four, proving the degree bound, including cases where it is identically zero.
With $T=\operatorname{diag}(1,-1)$, each matrix in \eqref{blocks} obeys
$M(-v)=TM(v)T$.  Replace $J(v)$ by
$\tfrac12\{J(v)+TJ(-v)T\}$; convexity of the positive semidefinite cone retains
both inequalities and gives the displayed parity.

On C1, $P(0)=\operatorname{diag}(0,4m+4\beta+2d)$ and
$N(0)=\operatorname{diag}(4C+4D+2d,0)$.  Zero diagonal entries in
$\mathcal G_J(0)$ force $p_0=s_0=0$. Finally
\[
q_1^2v^2=|J_{12}(v)|^2\leq4m(4D+2d)v^4
\]
for every $v$, so $q_1=0$. This proof does not assume any of the source weights
is positive.
\end{proof}

Use now the already established valid C1 source from
\texttt{strict\_outer\_physical\_source.json} in the public repository,
with $d=1$, $\lambda=n=0$, and
\begin{align*}
m&=\frac{1166523000}{31661820415017653},&
C&=\frac{322803000}{186199784281},&
A&=\frac{218288946200}{1359122513},\\
D&=\frac{90943434293}{222507771414},&
\alpha&=\frac{248802430236}{3300726103},&
\beta&=\frac{205476973}{51646655494}.
\end{align*}
Its validity is an established input from \cite{Project}, certified by the independently checked
order-13744 rational certificate; that certificate is not rebuilt here.
Its exact contact is
\[
u_*=\frac1{13},\qquad v_*=16,\qquad
(W_*,X_*,Y_*,Z_*)=(140,-1/32,1/50,1).
\]
In \eqref{change} the contact becomes
\begin{equation}
U_*=\frac1{169},\quad
a_* =\binom{69589/500}{247/5440},\qquad
b_* =\binom{69459/500}{-70473/5440}.\label{contact}
\end{equation}
The full four-dimensional equation
$\mathcal P_c(U_*,v_*)(a_*,b_*)^T=0$ is exact.

\begin{theorem}\label{thm:obstruction}
This valid source has no polynomial selector satisfying \eqref{lmis} for all
$v$. Equivalently, its polynomial $\Phi_c(u^2,v;a,b)$ is not a sum of polynomial
squares in $(u,v,a,b)$.
\end{theorem}
\begin{proof}
By Proposition~\ref{prop:degree}, a polynomial selector would have the form
\eqref{three}. At the contact, both nonnegative terms in \eqref{split} vanish.
Therefore
\begin{equation}
J(v_*)b_*=-P(v_*)a_*/U_*,\qquad
J(v_*)^Ta_*=-U_*N(v_*)b_*.\label{kernel}
\end{equation}
The first coordinate of the first equation determines $p$, and the second
coordinate of the second equation determines $s$; the remaining coordinates
then determine $r$ consistently. Substituting the displayed rational input gives
\begin{equation}
p=-\frac{37168625}{1489598274248},\qquad
r=\frac{622145600000}{505399414477},\qquad
s=\frac{7465747200}{3689046821}.\label{forced}
\end{equation}
This is a rank-three linear system, not a numerical fit.

Put $a_0=4m$, $b_0=4\beta+2d$, $c_0=4C$, $e_0=4D+2d$.  At $v=1$, the principal
submatrix of $\mathcal G_J$ on coordinates $(a_2,b_1,b_2)$ is
\[
\begin{pmatrix}
a_0+2b_0&r&s\\r&2c_0+e_0&e_0\\s&e_0&e_0
\end{pmatrix}.
\]
Its determinant equals
\begin{equation}
2e_0c_0(a_0+2b_0)-e_0(r-s)^2-2c_0s^2.\label{minor}
\end{equation}
The input fractions imply directly
\[
3<e_0<4,\quad 0<c_0<1/100,\quad 0<a_0+2b_0<5,\quad
r<5/4,\quad s>2.
\]
Consequently \eqref{minor} is less than
$2/5-27/16=-103/80<0$, contradicting $\mathcal G_J(1)\PSD$.

Finally, a polynomial sum of squares for $\Phi_c(u^2,v;a,b)$ can use roots
of degree at most two in $u$, at most one in $v$, and homogeneous linear
degree in the four state coordinates: higher terms cannot cancel in a sum
of real squares. Averaging in $u\mapsto-u$ and using the constant and leading
state supports exactly as in Theorem~\ref{thm:selector} would produce a
polynomial selector. Conversely a polynomial selector gives polynomial
matrix factorizations of its two positive univariate matrices by
Fej\'er--Riesz and then a polynomial sum of squares through \eqref{split}.
This proves the equivalence.
\end{proof}

\section{What remains, and the evaluated files}
The reduction is exact for the whole canonical slack cone, with both state
normals and every zero opening retained.  Its evaluated block diagonal
order-eight pencil $\mathcal G_J\oplus\mathcal H_J$ has 108 integer COO entries,
nine source parameters, four selector parameters, and Bob degree at most two.
The derived C1 three-parameter version has 80 integer COO entries.  They are
\emph{pointwise pencils}: the quantifier $\forall v\,\exists J$ remains.

The obstruction is to every polynomial Bob selector on the entire real Bob
chart in this normal form,
not to every source-independent algebraic family or multiplier.  It also
does not show that $u^2D_uF_c$ in the original state coordinates has no
polynomial sum-of-squares representation: cancelling the extra factors after
pullback through \eqref{change} would require another theorem.  A viable
completion must supply one fixed testing algebra (for example one fixed
Bob denominator or a genuinely algebraic family) in which the selectors and
their positive matrix factors exist uniformly over the full source cone.
The pointwise real-closed-field theorem and source-dependent matrix
factorization do not supply that assertion.  In particular, an odd algebraic
function such as $v|v|$ can vanish quadratically at $v=0$ without vanishing
identically, so the rank argument forcing the polynomial entry $J_{12}=0$
does not extend to a two-half-line algebraic cover.  No adequacy statement
for that cover has been proved here.

A concrete sufficient remaining implication is the existence of one fixed
nonnegative polynomial $q(v)$, positive on a dense set, for which every valid
source admits a polynomial $Z_c(v)$ satisfying
\[
\begin{pmatrix}q(P+K)&Z_c\\Z_c^T&q(N+K)\end{pmatrix}\PSD,\qquad
\begin{pmatrix}q(P+E)&q(K-E)-Z_c\\q(K-E)-Z_c^T&q(N+E)\end{pmatrix}\PSD
\quad(v\in\R).
\]
The first inequality would force $\deg Z_c\leq\deg q+2$; univariate
matrix factorization would then give one fixed testing family for
$q(v)u^2D_uF_c$, with source-linear coefficient input and an exact
validity converse by density. Theorem~\ref{thm:obstruction} rules out $q=1$
for this route. It neither proves nor disproves the displayed implication
for another $q$, or the broader possibility of a fixed algebraic family.
No full-set affine $Q_{222}$ lift is claimed here.

The two integer COO files accompanying this note contain respectively the
pointwise order-eight pencil and its derived C1 polynomial-selector version.
The counterexample JSON file contains the exact contact linear system, its
forced solution, and the negative minor.  An independent verifier checks every
exported matrix coefficient and the counterexample using only Python's
standard-library rational arithmetic.  A separate symbolic checker verifies
the full nine-source polynomial identities. The supplied builder recreates
only this new data; the README gives the exact commands.

\begin{thebibliography}{2}
\bibitem{Zalar} A. Zalar, \emph{Matrix Fej\'er--Riesz theorem with gaps},
Theorem 1.2, \href{https://arxiv.org/abs/1503.06034}{arXiv:1503.06034};
\href{https://zalara.github.io/Papers/MatrixFejerRieszArxivOriginal.pdf}{author's full text}.
\bibitem{Project} \raggedright A. Chaturvedi, \emph{q222-finite-sdp},
\href{https://github.com/anubhav70747/q222-finite-sdp/tree/5fc0a0487bb94d7bdfa9929edb5f4f6f3ff50240}{commit 5fc0a0487bb9}.
Established input: \texttt{extracted-lifts/outer-candidates/}
\texttt{strict\_outer\_physical\_source.json}; the exact file is bundled
with the new certificate data.
\end{thebibliography}
\end{document}
